\section{Introduction}
\label{sec:intro}

Benchmark scores increasingly influence which open-weight LLMs practitioners
select, which research directions are pursued, and which claims of progress
are accepted. Leaderboards compress each model's behavior on a benchmark
into a single number, and model comparisons are routinely made by comparing
those numbers across models.

A benchmark score, however, is not produced by the model alone. It is
produced by a measurement pipeline: a prompt template and instruction
phrasing, a number of few-shot exemplars, a chat template applied to the
rendered prompt, an inference backend and its version, a numerical precision
or quantization scheme, a decoding strategy with its seeds, and the hardware
the model runs on. Each of these is an operator choice that a reasonable
evaluation could set differently without changing the scientific question
being asked.

Prior work has documented that such choices matter. Prompt formatting alone
can spread performance by tens of accuracy points
\citep{sclar2024quantifying}, single-prompt evaluations can be brittle
enough that ranking claims do not survive a multi-prompt analysis
\citep{mizrahi2024state}, and adversarial prompt perturbations reliably
degrade LLM performance \citep{zhu2024promptrobust}. Evaluation-harness
developers have catalogued reproducibility hazards spanning backend,
precision, and prompt handling \citep{biderman2024lessons}, and statistical
treatments of evaluation as an experiment have argued for interval-based
reporting \citep{miller2024adding}.

Nevertheless, many published comparisons---particularly comparisons of
open-weight models assembled from heterogeneous runs---still rest on a
single configuration per model, and the interaction between configuration
sensitivity and \emph{model ranking} is rarely quantified directly. A score
change that leaves the ordering intact supports different conclusions than
one that swaps first and second place. Quantifying that distinction requires
per-item paired artifacts measured across multiple configuration dimensions
and models, which few evaluation setups publish alongside their scores.

This paper addresses that gap with a controlled empirical study built on a
configuration-aware evaluation harness. The harness records a full manifest
for every run---model and tokenizer identity, backend and version, chat
template, prompt variant, decoding parameters, seed, quantization, hardware,
package versions, and a SHA-256-derived configuration hash---writes per-item
correctness records, and registers every completed run in a hash-deduplicated
registry that binds the artifact to its configuration. On this infrastructure
we run a one-factor-at-a-time (OFAT) study: one control configuration per
model, with a single factor varied at a time across prompt variant, greedy
seed, inference backend, quantization level, and an explicitly sampled
decoding arm, on a frozen 800-item benchmark slice held fixed across all
runs. OFAT isolates each dimension's marginal effect relative to the frozen
baseline; it does not estimate interaction effects between dimensions, a
limitation we return to in \cref{sec:limitations}.
The measured results are: (i) prompt configuration is the dominant source
of score variability, moving a model's overall accuracy by up to
17.4\,\pp{}---about six times the 2.6\,\pp{} control-configuration gap between
the two strongest models, whose control Wilson intervals overlap; (ii) a
single legitimate prompt change (adding five few-shot exemplars) reverses
the Phi-3.5-mini / Qwen2.5-3B ordering with disjoint 95\% Wilson intervals
($\taub=0.33$ against the control ranking, one pairwise reversal), while
the other seven shared configurations preserve the control ordering; (iii)
quantization (INT8/INT4 via bitsandbytes) and backend (Hugging Face
\texttt{generate} vs.\ vLLM on identical rendered prompts) move scores by
at most 2.9\,\pp{} and 2.4\,\pp{}, and neither change is statistically
significant after Holm correction over the full 27-contrast family; (iv)
under greedy decoding, changing the RNG seed is an exact no-op (0 discordant
items out of 800 in every contrast), whereas temperature-0.7 sampling
changes 9.9--24.1\% of per-item outcomes while shifting net accuracy by at
most 3.5\,\pp{}.

This paper makes the following contributions:
\begin{enumerate}
  \item \textbf{A configuration-aware evaluation framework} that records the
  full evaluation context (manifest, configuration hash, git commit,
  hardware, package versions) for every run, documents incomparability
  explicitly, and registers results in a hash-deduplicated registry with
  per-item artifacts (\cref{sec:framework,sec:repro}).
  \item \textbf{A controlled OFAT sensitivity study} over five factor
  families on frozen official benchmark slices, with paired per-item
  statistics---Wilson intervals, exact McNemar tests with Holm correction
  over the 27-contrast family, paired bootstrap difference intervals, and
  supplementary Cohen's $h$ effect sizes---quantifying how much each dimension moves
  scores (\cref{sec:results,sec:sensitivity}).
  \item \textbf{A ranking-stability analysis} that separates score movement
  from ordering movement: Kendall's $\taub$ against the control ranking,
  pairwise reversal counts, per-configuration ranks, pairwise win rates over
  shared configurations, and a seeded bootstrap of ranking
  stability---including a measured, interval-supported ranking reversal
  under a prompt change (\cref{sec:ranking}).
  \item \textbf{An analysis of decoding stochasticity versus determinism}
  quantifying both the per-item churn and the net drift induced by
  temperature-0.7 sampling, and the exact seed-invariance of greedy
  decoding, in the same harness (\cref{sec:sampling}).
  \item \textbf{A post-hoc failure audit} categorizing incorrect and
  unparseable outputs across all 24,800 generated outputs, with
  representative examples and failure category analysis
  (\cref{sec:failures}).
  \item \textbf{A planned factorial experiment} (Prompt $\times$ Runtime
  $\times$ Seed) to analyze interaction effects that OFAT cannot recover,
  with committed configuration (\cref{sec:factorial-plan}).
  \item \textbf{Reproducibility artifacts} binding every reported number to
  its configuration hash, run manifest, git commit, and raw per-item
  records, with scripts that regenerate every table and figure in this
  paper from those artifacts (\cref{sec:repro}).
\end{enumerate}

We frame the measurement question formally in \cref{sec:framework}, describe
the experimental setup in \cref{sec:setup}, and present results in
\cref{sec:results}--\cref{sec:ranking}. \cref{sec:discussion} discusses
implications and the distinction between score and ranking instability; the
provisional reliability composite and exploratory diagnostics appear there.
\cref{sec:limitations} states limitations explicitly, organized as threats
to validity. We claim no results
beyond the measured models, tasks, configurations, and hardware documented
here.

